\chapter{Automatische Unit-Tests}
\label{ch:quality/unit-testing}

Unit-Tests\index{unit-test@Unit-Test} sind ein unverzichtbares Werkzeug,
um nachzuweisen und dich zu vergewissern, dass dein Code korrekt ist.
Dieses Kapitel zeigt, wie du mit \emph{Busted\index{busted@Busted}}
Tests für Luanti-Mods und -Spiele schreibst. Funktionen zu testen, die
Luanti-Funktionen aufrufen, ist ziemlich schwierig. Zum Glück haben wir
im vorigen Kapitel~(Kapitel~\ref{ch:quality/clean-arch}) besprochen, wie
du deinen Code so aufbaust, dass sich das vermeiden lässt.

\section{Busted installieren}\label{busted-installieren}

Zuerst installierst du LuaRocks:

\begin{itemize}
\tightlist
\item
  Windows: Folge der Installationsanleitung im LuaRocks-Wiki\footnote{\url{https://github.com/luarocks/luarocks/wiki/Installation-instructions-for-Windows}}
  (englisch).
\item
  Debian/Ubuntu: \texttt{sudo apt install luarocks}
\end{itemize}

Dann installierst du Busted global:

\begin{Shaded}
\begin{Highlighting}[]
\NormalTok{sudo luarocks install busted}
\end{Highlighting}
\end{Shaded}

Zum Schluss prüfst du, ob es installiert ist:

\begin{Shaded}
\begin{Highlighting}[]
\NormalTok{busted {-}{-}version}
\end{Highlighting}
\end{Shaded}

\section{Dein erster Test}\label{dein-erster-test}

Busted ist das führende Framework für Unit-Tests in Lua. Es sucht nach
Lua-Dateien, deren Name auf \texttt{\_\allowbreak{}spec} endet, und
führt sie in einer eigenständigen Lua-Umgebung aus.

\begin{Shaded}
\begin{Highlighting}[]
\NormalTok{mymod/}
\NormalTok{├── init.lua}
\NormalTok{├── api.lua}
\NormalTok{└── tests}
\NormalTok{    └── api\_spec.lua}
\end{Highlighting}
\end{Shaded}

\subsection{init.lua}\label{init.lua}

\begin{Shaded}
\begin{Highlighting}[]
\VariableTok{mymod} \OperatorTok{=} \OperatorTok{\{\}}

\FunctionTok{dofile}\OperatorTok{(}\VariableTok{core}\OperatorTok{.}\NormalTok{get\_modpath}\OperatorTok{(}\StringTok{"mymod"}\OperatorTok{)} \OperatorTok{..} \StringTok{"/api.lua"}\OperatorTok{)}
\end{Highlighting}
\end{Shaded}

\subsection{api.lua}\label{api.lua}

\begin{Shaded}
\begin{Highlighting}[]
\KeywordTok{function} \VariableTok{mymod}\OperatorTok{.}\NormalTok{add}\OperatorTok{(}\VariableTok{x}\OperatorTok{,} \VariableTok{y}\OperatorTok{)}
    \ControlFlowTok{return} \VariableTok{x} \OperatorTok{+} \VariableTok{y}
\KeywordTok{end}
\end{Highlighting}
\end{Shaded}

\subsection{tests/api\_spec.lua}\label{testsapi_spec.lua}

\begin{Shaded}
\begin{Highlighting}[]
\CommentTok{{-}{-} Benötigte Dateien im Mod{-}Verzeichnis suchen}
\VariableTok{package.path} \OperatorTok{=} \StringTok{"./?.lua;"} \OperatorTok{..} \VariableTok{package.path}

\CommentTok{{-}{-} Globale Tabelle mymod anlegen, in die die API schreibt}
\ConstantTok{\_G}\OperatorTok{.}\VariableTok{mymod} \OperatorTok{=} \OperatorTok{\{\}}
\CommentTok{{-}{-} api.lua ausführen}
\FunctionTok{require}\OperatorTok{(}\StringTok{"api"}\OperatorTok{)}

\CommentTok{{-}{-} Tests}
\NormalTok{describe}\OperatorTok{(}\StringTok{"add"}\OperatorTok{,} \KeywordTok{function}\OperatorTok{()}
\NormalTok{    it}\OperatorTok{(}\StringTok{"addiert"}\OperatorTok{,} \KeywordTok{function}\OperatorTok{()}
        \VariableTok{assert}\OperatorTok{.}\NormalTok{equals}\OperatorTok{(}\DecValTok{2}\OperatorTok{,} \VariableTok{mymod}\OperatorTok{.}\NormalTok{add}\OperatorTok{(}\DecValTok{1}\OperatorTok{,} \DecValTok{1}\OperatorTok{))}
    \KeywordTok{end}\OperatorTok{)}

\NormalTok{    it}\OperatorTok{(}\StringTok{"kann mit negativen Zahlen umgehen"}\OperatorTok{,} \KeywordTok{function}\OperatorTok{()}
        \VariableTok{assert}\OperatorTok{.}\NormalTok{equals}\OperatorTok{(}\DecValTok{0}\OperatorTok{,} \VariableTok{mymod}\OperatorTok{.}\NormalTok{add}\OperatorTok{({-}}\DecValTok{1}\OperatorTok{,} \DecValTok{1}\OperatorTok{))}
        \VariableTok{assert}\OperatorTok{.}\NormalTok{equals}\OperatorTok{({-}}\DecValTok{2}\OperatorTok{,} \VariableTok{mymod}\OperatorTok{.}\NormalTok{add}\OperatorTok{({-}}\DecValTok{1}\OperatorTok{,} \OperatorTok{{-}}\DecValTok{1}\OperatorTok{))}
    \KeywordTok{end}\OperatorTok{)}
\KeywordTok{end}\OperatorTok{)}
\end{Highlighting}
\end{Shaded}

Jetzt kannst du die Tests ausführen, indem du im Verzeichnis des Mods
ein Terminal öffnest und \texttt{busted .\allowbreak{}} aufrufst.

Wichtig ist, dass die API-Datei die Tabelle nicht selbst anlegt, denn
globale Variablen funktionieren in Busted anders. Jede Variable, die in
Luanti global wäre, ist in Busted lokal zur Datei. Das wäre auch für
Luanti der bessere Weg gewesen, aber dafür ist es jetzt zu spät.

Außerdem sollten die Dateien, die du testest, möglichst keine Funktionen
aufrufen, die nicht in ihnen selbst stehen. Meist schreibst du Tests
jeweils nur für eine einzelne Datei.

\section{Mocking: Externe Funktionen
ersetzen}\label{mocking-externe-funktionen-ersetzen}

\emph{Mocking\index{mocking@Mocking}} heißt, Funktionen zu ersetzen, von
denen das Getestete abhängt. Das hat zwei Zwecke: Zum einen steht die
Funktion in der Testumgebung womöglich nicht zur Verfügung, zum anderen
möchtest du vielleicht Aufrufe der Funktion und ihre Argumente
mitschneiden.

Hast du die Ratschläge aus dem Kapitel Saubere
Architektur~(Kapitel~\ref{ch:quality/clean-arch}) befolgt, hast du schon
eine recht saubere Datei zum Testen. Dinge außerhalb deines Bereichs
musst du aber trotzdem ersetzen -- etwa die View, wenn du Controller
oder API testest. Hast du die Ratschläge nicht befolgt, wird es etwas
schwieriger, weil du womöglich die Luanti-API ersetzen musst.

\begin{Shaded}
\begin{Highlighting}[]
\CommentTok{{-}{-} Wie oben eine Tabelle anlegen}
\ConstantTok{\_G}\OperatorTok{.}\VariableTok{core} \OperatorTok{=} \OperatorTok{\{\}}

\CommentTok{{-}{-} Die Mock{-}Funktion definieren}
\KeywordTok{local} \VariableTok{chat\_send\_player\_calls} \OperatorTok{=} \OperatorTok{\{\}}
\KeywordTok{function} \VariableTok{core}\OperatorTok{.}\NormalTok{chat\_send\_player}\OperatorTok{(}\VariableTok{name}\OperatorTok{,} \VariableTok{message}\OperatorTok{)}
    \FunctionTok{table.insert}\OperatorTok{(}
        \VariableTok{chat\_send\_player\_calls}\OperatorTok{,}
        \OperatorTok{\{} \VariableTok{name} \OperatorTok{=} \VariableTok{name}\OperatorTok{,} \VariableTok{message} \OperatorTok{=} \VariableTok{message} \OperatorTok{\}}
    \OperatorTok{)}
\KeywordTok{end}

\CommentTok{{-}{-} Tests}
\NormalTok{describe}\OperatorTok{(}\StringTok{"list\_areas"}\OperatorTok{,} \KeywordTok{function}\OperatorTok{()}
\NormalTok{    it}\OperatorTok{(}\StringTok{"gibt eine Zeile je Gebiet aus"}\OperatorTok{,} \KeywordTok{function}\OperatorTok{()}
        \VariableTok{chat\_send\_player\_calls} \OperatorTok{=} \OperatorTok{\{\}} \CommentTok{{-}{-} Tabelle zurücksetzen}

        \VariableTok{mymod}\OperatorTok{.}\NormalTok{list\_areas\_to\_chat}\OperatorTok{(}
            \StringTok{"singleplayer"}\OperatorTok{,}
            \StringTok{"singleplayer"}
        \OperatorTok{)}

        \VariableTok{assert}\OperatorTok{.}\NormalTok{equals}\OperatorTok{(}\DecValTok{2}\OperatorTok{,} \OperatorTok{\#}\VariableTok{chat\_send\_player\_calls}\OperatorTok{)}
    \KeywordTok{end}\OperatorTok{)}

\NormalTok{    it}\OperatorTok{(}\StringTok{"sendet an den richtigen Spieler"}\OperatorTok{,} \KeywordTok{function}\OperatorTok{()}
        \VariableTok{chat\_send\_player\_calls} \OperatorTok{=} \OperatorTok{\{\}} \CommentTok{{-}{-} Tabelle zurücksetzen}

        \VariableTok{mymod}\OperatorTok{.}\NormalTok{list\_areas\_to\_chat}\OperatorTok{(}
            \StringTok{"singleplayer"}\OperatorTok{,}
            \StringTok{"singleplayer"}
        \OperatorTok{)}

        \ControlFlowTok{for} \ConstantTok{\_}\OperatorTok{,} \VariableTok{call} \KeywordTok{in} \FunctionTok{pairs}\OperatorTok{(}\VariableTok{chat\_send\_player\_calls}\OperatorTok{)} \ControlFlowTok{do}
            \VariableTok{assert}\OperatorTok{.}\NormalTok{equals}\OperatorTok{(}\StringTok{"singleplayer"}\OperatorTok{,} \VariableTok{call}\OperatorTok{.}\VariableTok{name}\OperatorTok{)}
        \ControlFlowTok{end}
    \KeywordTok{end}\OperatorTok{)}

    \CommentTok{{-}{-} Die beiden Tests oben sind eigentlich überflüssig,}
    \CommentTok{{-}{-} denn dieser prüft beides}
\NormalTok{    it}\OperatorTok{(}\StringTok{"gibt das Richtige aus"}\OperatorTok{,} \KeywordTok{function}\OperatorTok{()}
        \VariableTok{chat\_send\_player\_calls} \OperatorTok{=} \OperatorTok{\{\}} \CommentTok{{-}{-} Tabelle zurücksetzen}

        \VariableTok{mymod}\OperatorTok{.}\NormalTok{list\_areas\_to\_chat}\OperatorTok{(}
            \StringTok{"singleplayer"}\OperatorTok{,}
            \StringTok{"singleplayer"}
        \OperatorTok{)}

        \KeywordTok{local} \VariableTok{expected} \OperatorTok{=} \OperatorTok{\{}
            \OperatorTok{\{}
                \VariableTok{name} \OperatorTok{=} \StringTok{"singleplayer"}\OperatorTok{,}
                \VariableTok{message} \OperatorTok{=} \StringTok{"Rathaus (2,43,63)"}\OperatorTok{,}
            \OperatorTok{\},}
            \OperatorTok{\{}
                \VariableTok{name} \OperatorTok{=} \StringTok{"singleplayer"}\OperatorTok{,}
                \VariableTok{message} \OperatorTok{=} \StringTok{"Flughafen (43,45,63)"}\OperatorTok{,}
            \OperatorTok{\},}
        \OperatorTok{\}}
        \VariableTok{assert}\OperatorTok{.}\NormalTok{same}\OperatorTok{(}\VariableTok{expected}\OperatorTok{,} \VariableTok{chat\_send\_player\_calls}\OperatorTok{)}
    \KeywordTok{end}\OperatorTok{)}
\KeywordTok{end}\OperatorTok{)}
\end{Highlighting}
\end{Shaded}

\section{Fazit}\label{fazit}

Gut eingesetzt, verbessern Unit-Tests die Qualität und Zuverlässigkeit
deines Projekts erheblich. Sie verlangen aber, dass du deinen Code
anders aufbaust als üblich.

Ein Beispiel für einen Mod mit vielen Unit-Tests ist crafting von
rubenwardy\footnote{\url{https://github.com/rubenwardy/crafting}}.
